\section{反常与 $\B_\cc$}\label{anombcc}

在 $\sigma$-模型微扰论的最低阶，辛流形 $Y$ 的每个拉格朗日子流形 $L$ 配上平坦 $\spinc$ 丛，都可伴随一个 $A$-膜 $\B_L$。然而在量子层面，盘瞬子效应会阻碍其中某些膜的存在。

对余各向同性膜的类似问题，物理学文献中显然尚未考虑，尽管一个与代数簇形变量子化相关的数学问题已被研究过 \cite{bezk}。秩 1 余各向同性膜由 Kapustin 与 Orlov 在 $\sigma$-模型微扰论最低阶构造出来 \cite{KO}。是否存在一个反常，会在更高阶阻碍某些此类膜的存在？

这里我们考虑支撑为整个 $A$-模型靶空间 $Y$ 的余各向同性膜 $\B_\cc$ 这一基本例子。特别地，$Y$ 是复辛流形，$\omega_Y=\Im\Omega$。复辛流形的第一 Chern 类为零，故通常的费米子数反常（它在分析盘瞬子对拉格朗日膜的阻碍时很重要）不存在，我们可以用费米子数守恒来约束分析。虽然我们只考虑支撑为整个 $Y$ 的余各向同性膜，但值得指出：至少对秩 1 的膜（这是余各向同性膜被理解得最好的唯一情形），一般情形是我们所考虑的情形与更熟悉的拉格朗日膜情形的混合。这是因为一般的秩 1 余各向同性膜在某些方向上像 $\B_\cc$，在另一些方向上像拉格朗日 $A$-膜。

若 $\B_\cc$ 是合法的 $A$-膜，这意味着理论的 BRST 算符 $Q$ 即使存在标记为 $\B_\cc$ 的边界时也满足 $Q^2=0$。另一种可能是 $Q^2$ 有边界贡献；虽然体内 $Q^2=0$，但可能有
\be\label{unu}Q^2=F(0),\ee
其中 $F(0)$ 是插在弦于 $\B_\cc$ 上结束的端点处的局部算符。弦的另一端点也可能有贡献，但我们可以只关注一个端点处的贡献。算符 $F(0)$ 在经典极限下共形维数为 0（若模型在量子层面不是共形不变的，这个标度可能被 $\sigma$-模型微扰论中的对数修正）。具有适当标度的算符形如 $F(X,\bar X,\rho)$，记号同附录 \ref{details}（这里场 $X,\bar X,\rho$ 在弦端点取值，即式 \eqref{unu} 中的点“0”）。如式 \eqref{btrans} 与 \eqref{bcorr} 的讨论所述，这些算符对应于 $Y$ 上的 $(0,q)$-形式，$q=0,\cdots,\dim_\C Y$，而 $Q$ 作用在它们上如同通常的 $\bar\partial$ 算符，把 $(0,q)$-形式映到 $(0,q+1)$-形式。$F(0)$ 费米子数为 2 的条件意味着它是 $(0,2)$-形式；条件 $Q^2=F(0)$ 蕴含 $[Q,F(0)]=[Q,Q^2]=0$，故 $F$ 是 $\bar\partial$-闭的 $(0,2)$-形式。若 $F$ 是 $\bar\partial$-恰当的，即对某个 $(0,1)$-形式 $\lambda$ 有 $F=\bar\partial\lambda$，则可通过平移 $Q\to Q-\lambda(0)$ 消除反常。因此，不能通过重新定义 $Q$ 消除的潜在反常，对应于 $H^2(Y,\O)$ 的元素。

然而，Bezrukavnikov 与 Kaledin 表明 \cite{bezk}（用不同的语言）：若存在 $Y$ 上的闭二形式 $b$，使其 $(0,2)$ 型部分（记为 $b^{(0,2)}$）等于 $\beta$，则与 $\beta\in H^2(Y,\O)$ 相伴的反常可以消除。为用我们的语言理解这一点，回忆所研究的 $A$-模型除辛形式 $\omega_Y=\Im\Omega$ 外，还有 B-场 $\sB=\Re\Omega$。作此特殊选择的理由是：Kapustin 与 Orlov \cite{KO} 在验证方程 $Q^2=0$ 没有边界贡献时，遇到条件 $(\omega_Y^{-1}(\sF+\sB))^2=-1$，它可由 $\sB=\Re\Omega$、$\sF=0$ 满足。若把 $\sB$ 平移为 $\sB+b$（$b$ 为闭二形式），则 Kapustin--Orlov 条件被破坏，$Q^2$ 得到正比于 $b^{(0,2)}$ 的边界贡献。但这也意味着，正比于 $b^{(0,2)}$ 的反常可通过平移 $\sB$ 抵消。

若 $Y$ 紧，则由 Hodge 理论，$H^2(Y,\O)$ 中每个类都有调和代表。这个代表是闭二形式，意味着任何潜在反常都总可通过平移 $\sB$ 抵消。通常对量子化并不关心 $Y$ 紧的情形，但上述论证对适当的一类非紧 $Y$ 大概有类比。例如，若 $Y$ 有渐近锥形的完备超 Kähler 度量，则可望证明任何反常都上同调于一个紧支撑 $(0,2)$-形式，然后用完备 Riemann 流形的 Hodge 理论说明反常可消去。不过我们不会沿此路线作任何精确论证。本文研究的大多数例子有 $H^{0,2}(Y)=0$。余切丛 $Y=T^*W$ 可能是例外。对这些例子，可用余切丛的伸缩对称证明没有反常。

Bezrukavnikov 与 Kaledin 用他们所谓的 period map 来参数化 $Y$ 上全纯函数层的可能形变量子化；这可看作 $\Omega$ 的上同调类 $[\Omega]$ 的量子版本。这样做时，会遇到 $Y$ 作为复辛流形可能有形变而不改变 $[\Omega]$ 这一点。此类形变的存在使 period map 的讨论复杂化。

从我们的观点看，$Y$ 作为复辛流形不改变 $[\Omega]$ 的形变有如下意义。设 $Y$ 的复结构与复辛形式 $I,\Omega$ 可形变为另一复结构与复辛形式 $I',\Omega'$，且 $[\Omega]=[\Omega']$。对 $(I',\Omega')$ 施加适当的微分同胚，可归到如下情形\footnote{一般地，若实辛形式 $[\omega]$ 在保持 $[\omega]$ 固定的连续族中变化，则该族中的辛形式在微分同胚意义下都等价。为看出这一点，注意保持 $[\omega]$ 固定的一阶变分形如 $\delta\omega=\d\gamma$（$\gamma$ 为某个 1-形式）。$\omega$ 的这一阶形变可由向量场 $V=\omega^{-1}\gamma$ 生成的无穷小微分同胚抵消。}：$\Im\Omega=\Im\Omega'$。不能同时要求 $\Re\Omega=\Re\Omega'$，但由于 $[\Re\Omega]=[\Re\Omega']$，确有 $\Re\Omega=\Re\Omega'-\d\lambda$（$\lambda$ 为某个 1-形式）。可把 $\sA=\lambda$ 视为平凡线丛 $\L$ 上的联络形式；$\L$ 的曲率为 $\sF=\d\lambda$。除熟悉的 Kapustin--Orlov 条件解 $(\omega_Y^{-1}(\sF+\sB))^2=-1$（$\omega_Y=\Im\Omega$、$\sB=\Re\Omega$、$\sF=0$）外，还有另一个解：$\omega_Y$ 与 $\sB$ 不变，但 $\sF=\d\lambda$（故 $\sF+\sB=\Re\Omega'$）。因此，除带平凡平坦 CP 丛的熟悉余各向同性膜 $\B_\cc$ 外，$\omega_Y=\Im\Omega$、$\sB=\Re\Omega$ 的 $A$-模型还有一个新的、同样好的余各向同性膜 $\B_\cc'$，其 CP 丛为 $\L$。更一般地，若导致膜 $\B_\cc$ 存在的复辛结构 $I,\Omega$ 可在保持 $[\Omega]$ 固定的族中变化，则固定 $\omega_Y$ 与 $\sB$ 的 $A$-模型有一族余各向同性膜，由同样的数据标记。（这是经典描述；在某些情形中，可通过平移 $\sB$ 抵消的反常可能使图像复杂化。）

保持 $[\Omega]$ 固定而形变 $I,\Omega$ 的可能性可如下描述。考虑层的正合列
\be\label{tolimo}0\to\C\to\O\to\O/\C\to0,\ee
其中 $\C$ 是局部截面为复常数的层。与此短正合列相伴，上同调中有正合列，其一部分为
\be\label{polimo}\cdots\to H^1(Y,\C)\xrightarrow{j}H^1(Y,\O)\to H^1(Y,\O/\C)\xrightarrow{\delta}H^2(Y,\C)\xrightarrow{j}H^2(Y,\O)\to\cdots.\ee
可把 $\O$ 解释为 $Y$ 上（全纯）Hamilton 函数的层，把 $\O/\C$ 解释为全纯函数模常数所得的层。模去常数对于分析 $Y$ 作为复辛流形的形变是合适的，因为与 Hamilton 函数 $f$ 相伴的 Hamilton 向量场 $\Omega^{-1}\d f$ 在 $f$ 平移常数下不变。相应地，$Y$ 作为复辛流形的一阶形变由 $H^1(Y,\O/\C)$ 分类。对与 $x\in H^1(Y,\O/\C)$ 相伴的形变，$[\Omega]$ 的相应形变为 $[\Omega]\to[\Omega]+\delta(x)$。因此，$Y$ 作为复辛流形不改变 $[\Omega]$ 的形变由 $\delta$ 的核分类。列 \eqref{polimo} 的正合性意味着 $\delta$ 的核等于商 $H^1(Y,\O)/j(H^1(Y,\C))$，它参数化 $\bar\partial$-闭的 $(0,1)$-形式模去形如 $c^{(0,1)}$（$c$ 为闭 1-形式）者。该商参数化在 $\omega_Y$ 与 $\sB$ 都不变的固定 $A$-模型中 $\B_\cc$ 作为余各向同性 $A$-膜的一阶形变族。
